\chapter{CONSIDERAÇÕES FINAIS E TRABALHOS FUTUROS}
\label{cap:conclusao}

Este capítulo encerra o documento de qualificação consolidando o percurso realizado até aqui. Retoma a questão e os objetivos de pesquisa, sintetiza o que já foi alcançado na proposição da Taxonomia Multidimensional de Integração da IAG no Audiovisual, explicitando as contribuições e as limitações da versão preliminar do artefato, e delineia os trabalhos futuros que conduzirão a pesquisa até a dissertação final.

\section{Retomada da Questão e dos Objetivos}

A questão que orienta este mestrado pergunta \textit{como desenvolver e validar uma taxonomia multidimensional que organize a integração da IAG no audiovisual não interativo, considerando suas dimensões técnicas, operacionais e éticas} (Seção~\ref{sec:intro:objetivo}). Diante da fragmentação da literatura e do descompasso entre o avanço técnico das ferramentas generativas e a sistematização de suas implicações sociotécnicas, o objetivo geral foi definido como desenvolver e validar um artefato classificatório capaz de correlacionar autonomia criativa, exigências técnicas e riscos éticos associados a cada escolha tecnológica.

Os três objetivos específicos (Seção~\ref{subsec:intro:objetivos_especificos}) situam-se, neste momento da pesquisa, em estágios distintos de cumprimento:

\begin{enumerate}
    \item \textbf{Sistematizar a base empírica.} Cumprido. O Capítulo~\ref{cap_trabalhos_previos} reapresenta o MSL de \citet{RothAleo2026} na forma em que foi originalmente relatado, submetendo o mesmo conjunto de evidências (sessenta estudos primários e dezessete subquestões, correspondentes a vinte e um campos de extração) a uma releitura analítica orientada pela meta-característica desta pesquisa. Onde o relato de origem organizava os dados para responder \textit{o que} a literatura reporta sobre IAG no audiovisual, aqui eles foram reexaminados para responder o que desse material efetivamente qualifica uma solução quanto à sua integração no \textit{pipeline} de produção não interativo. Desse reexame decorreram três operações analíticas que não existiam no estudo anterior: (i) a triagem das subquestões por poder discriminante, que manteve quinze delas como fundamento das facetas e excluiu deliberadamente SQ1 e SQ2 por descreverem propriedades do artigo, e não do objeto classificado (Tabela~\ref{tab:sq_dimensoes}); (ii) o reagrupamento das subquestões remanescentes em três eixos de mesma natureza analítica (operacional, técnica e sociotécnica), em substituição à apresentação subquestão por subquestão do relato original, tornando rastreável a origem empírica de cada dimensão do artefato; e (iii) a conversão das oito lacunas estruturais de achados descritivos em requisitos de projeto da taxonomia, com destaque para as lacunas de reprodutibilidade, de democratização e interdisciplinar, que se traduziram, respectivamente, nas facetas de transparência e reprodutibilidade, de infraestrutura e de impacto no trabalho. A base empírica é, portanto, herdada e devidamente atribuída; a interpretação que dela se faz neste documento é nova, subordinada à construção do artefato e sem correspondente no trabalho que a originou.

    \item \textbf{Propor a Taxonomia Multidimensional.} Cumprido em versão preliminar. O Capítulo~\ref{cap:taxonomia} apresentou o artefato facetado estruturado em três dimensões ortogonais: operacional (grau de autonomia, etapa do \textit{pipeline} e modo de interação), técnica (família arquitetural, modelo base, tipo de saída, infraestrutura e gargalo dominante) e de risco e ética (autoria, viés, impacto no trabalho, transparência/reprodutibilidade e perceptibilidade), acompanhado de síntese tabular, rubrica de decisão, protocolo de aplicação e leitura dos \textit{trade-offs} entre dimensões.

    \item \textbf{Validar a taxonomia proposta.} Em andamento. A validação plena (painel de especialistas e aplicação com medida de concordância entre avaliadores) constitui o objeto da continuidade desta pesquisa, conforme instrumentos já especificados na Seção~\ref{sec:tax:validacao}. Até o momento, realizou-se apenas uma verificação inicial de aplicabilidade sobre três estudos reais do corpus (Seção~\ref{sec:tax:instanciacao}), suficiente para demonstrar o poder descritivo do modelo e para refinar a rubrica, mas insuficiente para atestar validade científica e profissional definitivas.
\end{enumerate}

Em síntese, este documento de qualificação entrega a proposição fundamentada e operacionalizada do artefato; a etapa de validação e o ciclo iterativo de refinamento que dela decorrem permanecem como horizonte da dissertação.

\section{Síntese do Percurso Realizado}

O trabalho partiu do reconhecimento de que a IAG já atua de forma transversal no \textit{pipeline} audiovisual não interativo, da roteirização e do \textit{storyboard} à pós-produção e aos efeitos visuais, porém sob um conhecimento acadêmico disperso, com métricas heterogêneas e com taxonomias correlatas que isolam capacidades geradoras, riscos ou etapas produtivas. A fundamentação teórica (Capítulo~\ref{cap:conceitos}) delimitou o recorte no audiovisual linear, caracterizou as arquiteturas generativas relevantes e evidenciou a lacuna de um modelo que articule, simultaneamente, autonomia operacional, condicionantes técnicos e custos sociotécnicos.

Sobre essa lacuna, o MSL prévio forneceu a base empírica indispensável à abordagem indutiva de construção taxonômica \citep{nickerson2013method}. A análise dos sessenta estudos confirmou padrões recorrentes: hegemonia de Modelos de Difusão e \textit{Transformers}; concentração em poucos modelos-base comerciais; coerência temporal como gargalo técnico dominante; barreiras de infraestrutura computacional; e preocupações éticas centradas em impacto no trabalho, autoria e transparência. As lacunas de escala, reprodutibilidade, interdisciplinaridade e democratização, entre outras, tornaram explícita a necessidade de um artefato que não apenas catalogue ferramentas, mas explique relações e \textit{trade-offs}.

A partir das subquestões do MSL subordinadas à meta-característica, \textit{integração da IAG no pipeline de produção audiovisual não interativo}, a taxonomia proposta converteu evidências empíricas em dimensões, facetas e valores categóricos. Seu diferencial reside na natureza facetada e multidimensional: uma instância recebe, ao mesmo tempo, um perfil operacional, um perfil técnico e um perfil de risco. A rubrica de decisão transforma essa estrutura em método reutilizável, ao tornar explícitos os critérios de atribuição de cada valor. A aplicação a três estudos com graus de autonomia distintos (\textit{wr-AI-ter}, \textit{MagicVFX} e \textit{Anim-Director}) confirmou empiricamente a correlação entre maior autonomia delegada à máquina e maior custo sociotécnico, e já alimentou refinamentos concretos do artefato (valor \textit{texto/roteiro}, facetas multivaloradas e distinção entre risco \textit{baixo} e \textit{não reportado}).

\section{Principais Contribuições}

No estágio atual da pesquisa, as contribuições desta qualificação podem ser sintetizadas em quatro frentes:

\begin{itemize}
    \item \textbf{Contribuição teórico-metodológica.} A proposição de uma taxonomia multidimensional e facetada que integra, em um único artefato, as dimensões operacional, técnica e de risco/ética da adoção da IAG no audiovisual não interativo, preenchendo a lacuna deixada por modelos unidimensionais existentes na literatura.

    \item \textbf{Contribuição empírica de síntese.} A reorganização e a interpretação crítica dos resultados do MSL de \citet{RothAleo2026} sob a ótica da meta-característica taxonômica, tornando rastreável a origem de cada faceta nas subquestões de extração e evidenciando lacunas que orientam tanto a estrutura do artefato quanto a agenda de pesquisa futura.

    \item \textbf{Contribuição operacional.} A entrega de instrumentos concretos de uso do modelo (tabela-síntese, rubrica de decisão, protocolo em cinco passos e mapa de \textit{trade-offs}) que permitem classificar estudos e ferramentas de forma comparável, e não apenas descritiva.

    \item \textbf{Contribuição sociotécnica preliminar.} A explicitação de que ganhos de autonomia e eficiência técnica não são neutros: correlacionam-se com riscos de autoria, impacto no trabalho, opacidade de dados e perceptibilidade do material sintético, oferecendo um ``mapa de prontidão'' para decisões de adoção mais informadas.
\end{itemize}

Essas contribuições não encerram o ciclo da \textit{Design Science Research} \citep{hevner2004}: a avaliação rigorosa do artefato é condição para consolidá-las como resultado científico maduro.

\section{Limitações}

As limitações reconhecidas neste estágio são inseparáveis do caráter preliminar do artefato e do escopo próprio de um documento de qualificação:

\begin{itemize}
    \item A taxonomia ainda não passou por validação com painel de especialistas nem por aplicação ampliada com medida de concordância entre avaliadores; a utilidade e a reprodutibilidade do modelo permanecem, portanto, hipóteses fundamentadas, porém não plenamente testadas.

    \item A verificação de aplicabilidade cobriu apenas três estudos do corpus, classificados por um único codificador, o que restringe a generalização dos achados sobre o poder discriminante do artefato.

    \item A base empírica herda o recorte temporal e as fontes do MSL prévio (janela de busca de 2022 a abril de 2025 em IEEE Xplore e ACM Digital Library, com publicações efetivas entre 2023 e abril de 2025). Novos modelos, modalidades e regimes regulatórios podem exigir extensão de facetas, possibilidade prevista pelo critério de extensibilidade, mas ainda não exercitada.

    \item As escalas da dimensão de risco são qualitativas e dependem de julgamento informado pela rubrica; não pretendem mensurar impacto absoluto, e sim viabilizar comparação relativa entre soluções.

    \item Este documento não atualiza quantitativamente o MSL com nova rodada de busca; a opção metodológica, justificada na Seção~\ref{sec:intro:metodologia} pelo volume crescente de publicações e pela rápida obsolescência das ferramentas comerciais, privilegiou a síntese qualitativa e a construção do artefato sobre a base já consolidada, deixando eventuais ampliações longitudinais para etapas posteriores, se pertinentes à validação.
\end{itemize}

Tais limitações não invalidam a proposição; delimitam com precisão o que a continuidade da pesquisa precisa examinar e documentar.

\section{Trabalhos Futuros}

A continuidade desta pesquisa de mestrado organiza-se em torno da validação e do refinamento iterativo da taxonomia, em conformidade com o método de \citet{nickerson2013method} e com o plano já detalhado na Seção~\ref{sec:tax:validacao}. Os caminhos prioritários são:

\begin{enumerate}
    \item \textbf{Validação por painel de especialistas.} Submeter a versão preliminar a um painel de cinco a oito especialistas em IA aplicada e em produção audiovisual, dimensionamento justificado na Seção~\ref{sec:tax:validacao} pela necessidade de representar os dois perfis com redundância, de manter estável a mediana usada como critério de aceitação e de preservar a viabilidade da análise temática dos comentários abertos, utilizando o instrumento da Tabela~\ref{tab:instrumento_especialistas}, com critério de aceitação por mediana e análise temática dos comentários abertos para priorizar ajustes de clareza, completude, concisão, robustez, utilidade e extensibilidade.

    \item \textbf{Validação por aplicação e concordância entre avaliadores.} Classificar, de forma independente e com base exclusiva na rubrica, uma amostra estratificada de estudos do corpus acrescida de ferramentas comerciais vigentes; quantificar a concordância por faceta por meio do coeficiente \textit{kappa} de Cohen e revisar facetas que não atinjam ao menos concordância substancial.

    \item \textbf{Refinamento do artefato.} Incorporar os resultados das duas etapas de validação em ciclos documentados de desdobramento, fusão ou renomeamento de categorias, até a convergência entre aprovação do painel e reprodutibilidade da rubrica.

    \item \textbf{Ampliação do escopo de instanciação.} Explorar a capacidade do modelo para além do corpus científico, classificando \textit{pipelines} e ferramentas comerciais em uso na indústria, de modo a testar a utilidade prática do ``mapa de prontidão'' junto a profissionais criativos.

    \item \textbf{Desdobramentos de pesquisa.} Após a consolidação da versão validada, investigar usos derivados da taxonomia. Por exemplo, como estrutura de apoio à comparação de soluções, à identificação sistemática de lacunas ou à formulação de diretrizes de adoção responsável sem perder o foco na contribuição central do mestrado: o artefato taxonômico em si.
\end{enumerate}

\section{Considerações Finais}

Este documento de qualificação demonstrou que é possível, e necessário, organizar a integração da IAG no audiovisual não interativo por meio de um artefato que trate simultaneamente autonomia, técnica e risco. A Taxonomia Multidimensional proposta desloca o foco de ferramentas efêmeras para dimensões estruturantes, torna legíveis os \textit{trade-offs} da adoção tecnológica e oferece um protocolo objetivo para classificação reprodutível.

O caminho que resta (validar, refinar e consolidar) não é apêndice opcional, mas parte constitutiva do método adotado. Ao concluir esta etapa, a pesquisa encontra-se preparada para submeter o artefato ao escrutínio de especialistas e à prova empírica da aplicação independente, com o intuito de entregar, na dissertação, uma taxonomia madura, útil à academia e à indústria, e alinhada ao compromisso de uma adoção informada e responsável da IAG no audiovisual linear.
